\documentclass[a4paper, landscape]{article}

\input{include/cardLibs.tex}
\input{include/cardColours.tex}
\input{include/cardStyles.tex}

\begin{document}
\pagestyle{empty}

\begin{center}
\greycard{Ahab}{Ahoy, marujo!}{Você vence se Moby receber a condição de morto e você não.} 
\greycard{Moby}{Splash...}{Você vence se Ahab receber a condição de morto e você não.}
\greycard{Esposa}{Primeira-dama.}{Você vence se, ao fim do jogo, estiver na mesma sala que o Presidente e a Amante não estiver.} 
\greycard{Amante}{A outra mulher.}{Você vence se, ao fim do jogo, estiver na mesma sala que o Presidente e a Esposa não estiver.}
\greycard{Tocador}{Não consegue ficar quieto!}{Você vence se tocar em todos os jogadores até o fim do jogo.\\[3mm] Se o Puritano segurar seu pulso, você perde.} 
\greycard{Puritano}{Não solte!}{Você vence se segurar o pulso do Tocador. Você só pode segurar o pulso de uma pessoa por jogo.\\[3mm] Depois de segurar, não solte até o fim. Vocês não podem ser separados: ambos devem ser reféns, ou nenhum dos dois.}
\greycard{Julieta}{Encontre Romeu!}{Você vence se, ao fim do jogo, estiver na mesma sala que Romeu e ambos receberem a condição de morto.} 
\greycard{Romeu}{Encontre Julieta!}{Você vence se, ao fim do jogo, estiver na mesma sala que Julieta e ambos receberem a condição de morto.} 
\end{center}

\backpanel
\newpage

\begin{center}
\greycard{Atirador}{Um tiro. Um alvo.}{Ao fim do jogo, antes de qualquer pessoa revelar a carta, anuncie em público em quem você atira. Pode ser alguém de qualquer sala. Você vence se atirar no Alvo.\\[3mm] A pessoa atingida morre, ativando os substitutos disponíveis. Um Bombista morto não explode.} 
\greycard{Alvo}{Mantenha o segredo!}{Você vence se o Atirador não atirar em você ao fim do jogo.} 
\greycard{Isca}{Chame a atenção!}{Você vence se o Atirador atirar em você ao fim do jogo.} 
\greycard{Pedra}{Firme como uma rocha!}{Você vence se, ao fim do jogo, estiver na mesma sala que Tesoura e Papel não estiver.} 
\greycard{Papel}{Papel em ação!}{Você vence se, ao fim do jogo, estiver na mesma sala que Pedra e Tesoura não estiver.} 
\greycard{Tesoura}{Bem afiada!}{Você vence se, ao fim do jogo, estiver na mesma sala que Papel e Pedra não estiver.}
\greycard{Mordomo}{Às suas ordens.}{Você vence se, ao fim do jogo, estiver na mesma sala que a Empregada e o Presidente.} 
\greycard{Empregada}{Às suas ordens.}{Você vence se, ao fim do jogo, estiver na mesma sala que o Mordomo e o Presidente.} 
\end{center}

\backpanel
\newpage

\begin{center}
\greycard{Rival}{Meu inimigo!}{Você vence se não estiver na mesma sala que o Presidente ao fim do jogo.} 
\greycard{Sobrevivente}{Não vou desistir!}{Você vence se não receber a condição de morto.} 
\greycard{Estagiário}{Dê o seu melhor!}{Você vence se estiver na mesma sala que o Presidente ao fim do jogo.} 
\greycard{Anarquista}{Que o mundo pegue fogo!}{Você vence se ajudar a destituir um Líder na maioria das rodadas.}
\greycard{Lacaio}{O mestre deve liderar!}{Você vence se nenhum Líder for destituído na mesma sala que você durante o jogo.}
\greycard{Vítima}{Essa bomba é para você!}{Você vence se receber a condição de morto devido à explosão do Bombista.} 
\end{center}

\backpanel

\end{document}